\documentclass[a4paper, 12pt]{article}

\usepackage{utils}

\renewcommand*{\today}{16 January 2026}

\begin{document}

\hotbox{Analyse 4}{CM 3}{\today}

\begin{proposition}{Dérivabilité}{}
    Soit $I$ un intervalle de $\R$, $(f_n)$ une suite de fonctions de $I$ dans $\R$ et $g: I \to \R$.

    \begin{itemize}[label=--]
        \item Pour tout $n \in \N$, $f_n$ est de classe $C^1$ sur $I$
        \item Il existe un point $x_0 \in I$ tel que $(f_n(x_0))$ converge
        \item $(f_n')$ converge uniformément vers $g$ sur tout compact de $I$.
    \end{itemize}
    Alors $(f_n)$ converge uniformément vers $f$ sur tout compact $J \subset I$, \text{$f$ est $C^1$ et $f' = g$}.
\end{proposition}

\begin{demonstration}
    Soit $y \in I$, on peut construire $J = [a, b] \subset I$ avec $a = \min(x_0, y)$ et $b = \max(x_0, y)$.

    On pose $\funcdef{G_n}{I}{\R}{x}{\int_{x_0}^x f_n'(t)dt}$, $\funcdef{G}{I}{\R}{x}{\int_{x_0}^x g(t)dt}$ et $l := \lim\limits_{n \to +\infty} f_n(x_0)$.

    En appliquant le théorème \ref{theoreme:integrale} à $(f_n')$, on obtient que $G_n$ converge uniformément vers $G$ sur $[a, b]$.

    Par ailleurs, d'après le théorème fondamental de l'analyse, on a $G_n(x) = f_n(x) - f_n(x_0)$ et donc $f_n(x) = G_n(x) + f_n(x_0)$.
    
    On note $\funcdef{f}{I}{\R}{x}{G(x) + l}$ par passage à la limite ($y$ étant arbitraire dans $I$).

    $g$ est continue comme limite uniforme de fonctions continues $(f_n')$ sur tout compact de $I$ donc $G$ est de classe $C^1$ sur $I$ donc $f$ aussi (somme de fonctions $C^1$).

    De plus $f'(x) = G'(x) = g(x)$.
    
    \vspace{0.5cm}

    Il reste à montrer la convergence uniforme de $(f_n)$ vers $f$ sur $J$.

    Ainsi pour tout $x \in J$, on a
    \begin{align*}
        |f_n(x) - f(x)| &= |G_n(x) + f_n(x_0) - G(x) - l| \\
        &\leq |G_n(x) - G(x)| + |f_n(x_0) - l| \\
        &= |\int_{x_0}^x (f_n'(t) - g(t)) dt| + |f_n(x_0) - l| \\
        &\leq |x - x_0| \|f_n' - g\|_{\infty, J} + |f_n(x_0) - l| \\
        &\leq |b - a| \|f_n' - g\|_{\infty, J} + |f_n(x_0) - l|
    \end{align*}

    Donc, $\sup\limits_{x \in J} |f_n(x) - f(x)| \leq |b - a| \|f_n' - g\|_{\infty, J} + |f_n(x_0) - l|$.

    Le second membre ne dépend pas de $x$ et par hypothèse tend vers $0$ quand $n$ tend vers $+\infty$.

    On en déduit que $(f_n)$ converge uniformément vers $f$ sur $J$ et comme $y$ était arbitraire dans $I$, la convergence est uniforme sur tout compact de $I$.

\end{demonstration}

\end{document}